The case was assembled measurement by measurement over four decades. The Alvarez team reported an anomalous enrichment of iridium in the Cretaceous–Paleogene boundary clay at Gubbio, Italy. Iridium is scarce in Earth's crust and comparatively abundant in meteoritic material. The Gubbio measurements came first; only afterwards was the same anomaly documented at more than 100 Cretaceous–Paleogene boundary sections worldwide. The Alvarez iridium results were published in the journal Nature in 1980 — that is, the enrichment measured by that team in the Italian boundary clay appeared as a 1980 Nature paper.

Much of the dispute turned on what the boundary layer itself contains, and on how the minerals in it could have been made. Shocked quartz grains are present in the Cretaceous–Paleogene boundary layer, and the petrography of those grains was examined at length by partisans of both the impact and the volcanic hypotheses.

Glen Penfield and Antonio Camargo identified a buried circular structure beneath the Yucatán Peninsula during petroleum surveys for Pemex. More than a decade after that identification, Alan Hildebrand and colleagues published a 1991 paper linking the Chicxulub structure to the Cretaceous–Paleogene boundary. The Deccan Traps form a flood basalt province covering much of west-central India. The shocked quartz at the boundary layer formed during explosive eruptions of the Deccan Traps, a derivation that cannot be squared with the standing of the Deccan Traps as a flood basalt province covering much of west-central India. The eruptive origin of those grains is the evidence on which Gerta Keller argued that Deccan Traps volcanism was the primary driver of the end-Cretaceous extinction. The correlation established in the 1991 Hildebrand paper indicates, for its part, that the Chicxulub impact and the Cretaceous–Paleogene boundary are coeval within radiometric dating uncertainty. Either that indistinguishability of ages holds, or the Chicxulub impact predates the Cretaceous–Paleogene boundary by roughly 300,000 years: the crater is either synchronous with the boundary or some 300,000 years older than it, and it cannot be both.

\slotmarker{S1} Deep drilling into the crater structure thus supplied the physical material against which the competing accounts of the extinction could be tested, and it is on such direct sampling, rather than on argument from correlation alone, that the judgement of the geological community finally rested.
